\section{The field's considered rejection of attacker compute}
\label{sec:findings}

Part of the field has considered pricing the attacker and declined. In a set
of recommendations for evaluating adversarial robustness, circulated as a
preprint and widely adopted without formal publication, Carlini et al.\ write
that there ``are few reasonable threat models under which an attacker can
compute 100 iterations of gradient descent, but not 1000''. They conclude that
an evaluation should assume an adversary who spends the time the attack needs
\cite{CAPBRTGMK19}. Their previous sentence carries the principle, ``Because the threat model must
not constrain the approach an attacker takes''. This is a considered position
and not an oversight, and it matches the practice of earlier work by an
overlapping set of authors, who run as many iterations of gradient descent as
convergence needs \cite{ACW18}.
% claim: field-rejects-attacker-compute-as-a-parameter

The position rejects attacker compute as a restriction imposed inside a fixed
threat model. It does not address attacker compute as the unit in which two
claims are compared, and work in the same area adopts that second use. Suya et
al.\ survey \BBattacks{} black-box attacks, re-run \BBrerun{} transfer attacks
under a fixed budget of \BBminutes{} minutes per batch, and recommend that
``execution cost such as the local attack runtime should be used as the
equalizing factor when comparing black-box attack performance, not the number
of iterations'' \cite{SSZHTE24}. Their runtime enters no definition of
robustness, so a defence still earns no credit for costing the attacker more.
